\scanpage{204}

i.e. qui sont loc. des dérivées logarithmiques, permettent de construire des \deleted{\unsure{valeurs|vecteurs}}\inserted{\unsure{vecteurs|valeurs}} propres de $F_{\deleted{X/S}}^{\nu}$ \deleted{qui se \unsure{x} de \unsure{$\omega$}} ($\ell$-adique, ou cristallin) qui sont $\equiv1$ $(p)$.

Plus précisément, il existe $\omega$ avec $\widetilde C\omega=\widetilde\omega$, $\omega\ne0$, ssi il existe une valeur propre de Frobenius, congrue à $1$ mod $p$. La dim de l'espace \deleted{des} engendré par les $\omega$ (\uuline{\unsure{vect.|vectr.} sur} $k$ !) tels que $\widetilde C\omega=\widetilde\omega$ est égale au nb des valeurs propres de Frobenius congrues à $1$ mod $p$. C'est aussi égal à la dim sur $\mathbb F_p$ du groupe \deleted{des \unsure{pts|pt} de la jacobienne} $ {}_pA(k)$, $A$ la jacobienne.

\margin{Fin 2\textsuperscript{e} partie ..}

% archive sheet number: 203
